% !TEX program = pdflatex
% Example: Elsevier journal article template
% Compile with: pdflatex main.tex && bibtex main && pdflatex main.tex && pdflatex main.tex
% Requires: elsarticle class (tlmgr install elsarticle)

\documentclass[review,3p,times]{elsarticle}

\usepackage{amsmath,amssymb}
\usepackage{graphicx}
\usepackage{booktabs}

\journal{Pattern Recognition}

\begin{document}

\begin{frontmatter}
\title{A Comparative Study of Feature Extraction Methods
       for Biomedical Signal Classification}

\author[a]{Author One}
\author[b]{Author Two}

\affiliation[a]{organization={Department of Biomedical Engineering, University A},
                addressline={123 University Street}, city={City}, country={Country}}
\affiliation[b]{organization={Department of Computer Science, University B},
                addressline={456 College Avenue}, city={City}, country={Country}}

\begin{abstract}
This paper presents a comparative study of feature extraction
methods for biomedical signal classification. We evaluate
time-domain, frequency-domain, and time-frequency approaches
on three public datasets.
\end{abstract}

\begin{keyword}
Biomedical signals \sep Feature extraction \sep Classification \sep Deep learning
\end{keyword}
\end{frontmatter}

\section{Introduction}
Biomedical signal classification is a fundamental task in
healthcare diagnostics. The choice of feature extraction
method significantly impacts classification performance.

\section{Methods}
We compare three categories of feature extraction:
\begin{itemize}
\item Time-domain: statistical moments, Hjorth parameters.
\item Frequency-domain: power spectral density, band powers.
\item Time-frequency: wavelet transform, Hilbert-Huang transform.
\end{itemize}

\section{Results}
Our experiments show that wavelet-based features achieve
the highest classification accuracy across all datasets.

\section{Conclusion}
Time-frequency methods consistently outperform time-domain
and frequency-domain approaches for biomedical signal classification.

\bibliographystyle{elsarticle-num}
\bibliography{references}

\end{document}
